% !TEX root = main.tex

% Theorem System
% The following boxes are provided:
%   Definition:     \defn 
%   Theorem:        \thm 
%   Lemma:          \lem
%   Corollary:      \cor
%   Proposition:    \prop   
%   Claim:          \clm
%   Fact:           \fact
%   Proof:          \pf
%   Example:        \ex
%   Remark:         \rmk (sentence), \rmkb (block)
% Suffix
%   r:              Allow Theorem/Definition to be referenced, e.g. thmr
%   p:              Add a short proof block for Lemma, Corollary, Proposition or Claim, e.g. lemp
%                   For theorems, use \pf for proof blocks

\tcbset{
    kb/common block/.style={
        enhanced,
        breakable,
        frame hidden,
        boxrule=0pt,
        sharp corners,
        boxsep=0pt,
        left=6mm,
        right=5mm,
        top=4mm,
        bottom=4mm,
        before skip=10pt,
        after skip=12pt,
        before upper={\parindent15pt\relax},
    },
    kb/proof block/.style={
        kb/common block,
        colback=black!1!white,
        borderline west={0.9mm}{0pt}{black!55},
    },
    kb/accent proof/.style n args={2}{
        kb/common block,
        colback=#1!4!white,
        borderline west={1mm}{0pt}{#1},
        before upper={
            \parindent15pt\relax
            {\noindent\bfseries #2}\par\smallskip
        },
    },
    kb/example block/.style={
        kb/common block,
        colback=MidnightBlue!3!white,
        borderline west={1.2mm}{0pt}{MidnightBlue!70!black},
        borderline north={0.3mm}{0pt}{MidnightBlue!18!white},
    },
    kb/example title/.style={
        colback=MidnightBlue!75!black,
        colframe=MidnightBlue!75!black,
        boxrule=0pt,
        sharp corners,
        left=3mm,
        right=3mm,
        top=1.2mm,
        bottom=1.2mm,
    },
}

\newcounter{example}[section]
\renewcommand{\theexample}{\thesection.\arabic{example}}

% Definition
\newtcbtheorem[number within=section]{mydefinition}{Definition}
{
    enhanced,
    frame hidden,
    titlerule=0mm,
    toptitle=1mm,
    bottomtitle=1mm,
    fonttitle=\bfseries\large,
    coltitle=black,
    colbacktitle=blue!20!white,
    colback=blue!10!white,
}{defn}

\NewDocumentCommand{\defn}{m+m}{
    \begin{mydefinition}{#1}{}
        #2
    \end{mydefinition}
}

\NewDocumentCommand{\defnr}{mm+m}{
    \begin{mydefinition}{#1}{#2}
        #3
    \end{mydefinition}
}

\NewDocumentEnvironment{definition}{o+b}{
    \IfNoValueTF{#1}{
        \begin{mydefinition}{}{}
    }{
        \begin{mydefinition}{#1}{}
    }
    #2
    \end{mydefinition}
}{}

% Theorem
\newtcbtheorem[use counter from=mydefinition]{mytheorem}{Theorem}
{
    enhanced,
    frame hidden,
    titlerule=0mm,
    toptitle=1mm,
    bottomtitle=1mm,
    fonttitle=\bfseries\large,
    coltitle=black,
    colbacktitle=cyan!20!white,
    colback=cyan!10!white,
}{thm}

\NewDocumentCommand{\thm}{m+m}{
    \begin{mytheorem}{#1}{}
        #2
    \end{mytheorem}
}

\NewDocumentCommand{\thmr}{mm+m}{
    \begin{mytheorem}{#1}{#2}
        #3
    \end{mytheorem}
}

\NewDocumentEnvironment{theorem}{o+b}{
    \IfNoValueTF{#1}{
        \begin{mytheorem}{}{}
    }{
        \begin{mytheorem}{#1}{}
    }
    #2
    \end{mytheorem}
}{}

% Lemma
\newtcbtheorem[use counter from=mydefinition]{mylemma}{Lemma}
{
    enhanced,
    frame hidden,
    titlerule=0mm,
    toptitle=1mm,
    bottomtitle=1mm,
    fonttitle=\bfseries\large,
    coltitle=black,
    colbacktitle=violet!20!white,
    colback=violet!10!white,
}{lem}

\NewDocumentCommand{\lem}{m+m}{
    \begin{mylemma}{#1}{}
        #2
    \end{mylemma}
}

\newenvironment{lempf}{
    \begin{tcolorbox}[kb/accent proof={violet!65!black}{Proof for Lemma.}]
}{
    \hfill\textcolor{violet!65!black}{$\blacksquare$}
    \end{tcolorbox}
}

\NewDocumentCommand{\lemp}{m+m+m}{
    \begin{mylemma}{#1}{}
        #2
    \end{mylemma}

    \begin{lempf}
        #3
    \end{lempf}
}

% Corollary
\newtcbtheorem[use counter from=mydefinition]{mycorollary}{Corollary}
{
    enhanced,
    frame hidden,
    titlerule=0mm,
    toptitle=1mm,
    bottomtitle=1mm,
    fonttitle=\bfseries\large,
    coltitle=black,
    colbacktitle=orange!20!white,
    colback=orange!10!white,
}{cor}

\NewDocumentCommand{\cor}{+m}{
    \begin{mycorollary}{}{}
        #1
    \end{mycorollary}
}

\newenvironment{corpf}{
    \begin{tcolorbox}[kb/accent proof={orange!80!black}{Proof for Corollary.}]
}{
    \hfill\textcolor{orange!80!black}{$\blacksquare$}
    \end{tcolorbox}
}

\NewDocumentCommand{\corp}{m+m+m}{
    \begin{mycorollary}{}{}
        #1
    \end{mycorollary}

    \begin{corpf}
        #2
    \end{corpf}
}

% Proposition
\newtcbtheorem[use counter from=mydefinition]{myproposition}{Proposition}
{
    enhanced,
    frame hidden,
    titlerule=0mm,
    toptitle=1mm,
    bottomtitle=1mm,
    fonttitle=\bfseries\large,
    coltitle=black,
    colbacktitle=yellow!30!white,
    colback=yellow!20!white,
}{prop}

\NewDocumentCommand{\prop}{+m}{
    \begin{myproposition}{}{}
        #1
    \end{myproposition}
}

\newenvironment{proppf}{
    \begin{tcolorbox}[kb/accent proof={Goldenrod!85!black}{Proof for Proposition.}]
}{
    \hfill\textcolor{Goldenrod!85!black}{$\blacksquare$}
    \end{tcolorbox}
}

\NewDocumentCommand{\propp}{+m+m}{
    \begin{myproposition}{}{}
        #1
    \end{myproposition}

    \begin{proppf}
        #2
    \end{proppf}
}

% Claim
\newtcbtheorem[use counter from=mydefinition]{myclaim}{Claim}
{
    enhanced,
    frame hidden,
    titlerule=0mm,
    toptitle=1mm,
    bottomtitle=1mm,
    fonttitle=\bfseries\large,
    coltitle=black,
    colbacktitle=pink!30!white,
    colback=pink!20!white,
}{clm}


\NewDocumentCommand{\clm}{m+m}{
    \begin{myclaim*}{#1}{}
        #2
    \end{myclaim*}
}

\newenvironment{clmpf}{
    \begin{tcolorbox}[kb/accent proof={Magenta!70!black}{Proof for Claim.}]
}{
    \hfill\textcolor{Magenta!70!black}{$\blacksquare$}
    \end{tcolorbox}
}

\NewDocumentCommand{\clmp}{m+m+m}{
    \begin{myclaim*}{#1}{}
        #2
    \end{myclaim*}

    \begin{clmpf}
        #3
    \end{clmpf}
}

% Fact
\newtcbtheorem[use counter from=mydefinition]{myfact}{Fact}
{
    enhanced,
    frame hidden,
    titlerule=0mm,
    toptitle=1mm,
    bottomtitle=1mm,
    fonttitle=\bfseries\large,
    coltitle=black,
    colbacktitle=purple!20!white,
    colback=purple!10!white,
}{fact}

\NewDocumentCommand{\fact}{+m}{
    \begin{myfact}{}{}
        #1
    \end{myfact}
}


% Proof
\RenewDocumentEnvironment{proof}{O{\proofname}}{
    \par
    \pushQED{\qed}
    \begin{tcolorbox}[kb/proof block]
        {\noindent\bfseries #1}\par\smallskip
        \ignorespaces
}{
    \popQED
    \end{tcolorbox}
    \par
}

\let\proofenvironment\proof
\let\endproofenvironment\endproof

\makeatletter
\renewcommand{\proof}{
    \@ifnextchar\bgroup{\proofcommand}{\proofenvironment}
}
\makeatother

\NewDocumentCommand{\proofcommand}{+m}{
    \proofenvironment
        #1
    \endproofenvironment
}

\NewDocumentCommand{\pf}{o+m}{
    \IfNoValueTF{#1}{
        \begin{proof}
            #2
        \end{proof}
    }{
        \begin{proof}[#1]
            #2
        \end{proof}
    }
}

% Example
\NewDocumentEnvironment{example}{o+b}{
    \refstepcounter{example}
    \IfNoValueTF{#1}{
        \def\exampletitle{Example \theexample}
    }{
        \def\exampletitle{Example \theexample \textbar\ #1}
    }
    \begin{tcolorbox}[
        kb/example block,
        title={\exampletitle},
        fonttitle=\bfseries\small,
        coltitle=white,
        boxed title style={kb/example title},
        attach boxed title to top left={xshift=4mm,yshift*=-\tcboxedtitleheight/2},
        top=7mm,
    ]
        #2
    \end{tcolorbox}
}{}

\NewDocumentCommand{\ex}{+m}{
    \begin{example}
        #1
    \end{example}
}


% Remark
\NewDocumentCommand{\rmk}{+m}{
    {\it \color{blue!50!white}#1}
}

\newenvironment{remark}{
    \par
    \vspace{5pt}
    \begin{minipage}{\textwidth}
        {\par\noindent{\textbf{Remark.}}}
        \tcolorbox[blanker,breakable,left=5mm,
        before skip=10pt,after skip=10pt,
        borderline west={1mm}{0pt}{cyan!10!white}]
}{
        \endtcolorbox
    \end{minipage}
    \vspace{5pt}
}

\NewDocumentCommand{\rmkb}{+m}{
    \begin{remark}
        #1
    \end{remark}
}
